% !TeX root = Foundations_of_Reinforcement_Learning.tex
% !TeX root = ./Roboter.tex


\usepackage[ngerman]{babel}
\usepackage[utf8]{inputenc}
\setlength{\parindent}{0em}
\usepackage{eurosym}
\usepackage{amsmath}
\usepackage{amssymb}
\usepackage{dsfont}
\usepackage{polynom}
\usepackage{graphicx}
\usepackage{caption}
\usepackage{tikz}
\usepackage[a4paper,portrait,left=1.0cm,right=1.0cm,top=2cm,bottom=2cm]{geometry}
\usepackage{pgf} % Zur Einbindung der PGF Files
\usepackage{bm}

\usepackage{footnote}


\usepackage{algorithm}
\usepackage{algpseudocode}  % pseudocode-Umgebung (\begin{algorithmic}, \State, \For, ...)


% Light Mode Customization
\usepackage{xcolor}
\pagecolor{white}  % Hintergrund auf Weiß
\color{black}      % Text auf Schwarz

% Eigene Farben definiert (nun optimiert für hellen Hintergrund)
\definecolor{rot}{RGB}{204,0,0}
\definecolor{blau}{RGB}{0,102,204}
\definecolor{orange}{RGB}{204,85,0}
\definecolor{gruen}{RGB}{0,153,0}
\definecolor{dunkelgruen}{RGB}{0,153,0}
\definecolor{dunkelblau}{RGB}{0,128,128}
\definecolor{backcolor}{RGB}{235,245,255}
\definecolor{frontcolor}{RGB}{0,153,0}
\definecolor{hellblau}{RGB}{180,210,235}
\definecolor{lila}{RGB}{153,0,153}
\definecolor{gelb}{RGB}{204,153,0}
\definecolor{dunkelrot}{RGB}{153,0,51}
\definecolor{dunkelgelb}{RGB}{153,153,0}
\definecolor{boxgray}{RGB}{245,245,245}        % Sehr helles Grau für Box-Hintergründe
\definecolor{boxframe}{RGB}{200,200,200}       % Helles Grau für Box-Rahmen

% Section colors (Headings)
%\usepackage{sectsty}
%\sectionfont{\color{textwhite}}
%\subsectionfont{\color{textwhite}}
%\subsubsectionfont{\color{textwhite}}

% Box customization mit mdframed
\usepackage{mdframed}
\mdfdefinestyle{mystyle}{
    leftmargin=0pt,
    rightline=true,
    topline=false,
    bottomline=false,
    leftline=false,
    linecolor=boxframe,
    linewidth=1pt,
    backgroundcolor=white,
    fontcolor=black,
    %frametitlefont={\color{gruen}\bfseries}
}

% Custom theorem environments
\newmdtheoremenv[
    style=mystyle,
    leftmargin=-10pt,
    linecolor=rot
    ]{satz}{Satz}[section] % linecolor=rot \color{rot}Satz
\newmdtheoremenv[
    style=mystyle,
    leftmargin=-10pt,
    linecolor=gelb
    ]{definition}[satz]{Definition} % \color{gelb}Definition
\newmdtheoremenv[
    style=mystyle,
    leftmargin=-10pt,
    linecolor=lila
    ]{beispiel}[satz]{Beispiel} % \color{lila}Beispiel
\newenvironment{beweis}{
    \begin{mdframed}[
            style=mystyle,
            leftmargin=-10pt,
            linecolor=gruen
        ]
        \textbf{Beweis:} % \color{gruen}Beweis:
        }{
        \hfill $\Box$
    \end{mdframed}
}

%\makesavenotesenv{beweis}
%\makesavenotesenv{beispiel}
%\makesavenotesenv{definition}
%\makesavenotesenv{satz}

% Listings für Python
%\usepackage[german,tworuled]{algorithm2e}
\usepackage{listings}
\lstset{
    commentstyle=\color{gruen},
    keywordstyle=\color{blau},
    language=Python,
    backgroundcolor=\color{boxgray},
    basicstyle=\ttfamily\color{black},
    stringstyle=\color{rot},
    tabsize=2,
    literate={@}{{@}}1,
    title=\lstname
}

% Hyperref Settings
\usepackage[pdfstartview=Fit,plainpages=false,colorlinks=true,linkcolor=gruen,linktocpage=true,hyperfootnotes=false,bookmarksopen=false]{hyperref}

% Eigene Kommandos
\newcommand{\e}[1]{\text{e}^{#1}}
\newcommand{\rot}[2][rot]{\textcolor{#1}{#2}}
\newcommand{\gruen}[2][gruen]{\textcolor{#1}{#2}}
\newcommand{\frontcolor}[2][frontcolor]{\textcolor{#1}{#2}}
\newcommand{\gelb}[2][gelb]{\textcolor{#1}{#2}}
\newcommand{\orange}[2][orange]{\textcolor{#1}{#2}}
\newcommand{\blau}[2][blau]{\textcolor{#1}{#2}}
\newcommand{\hellblau}[2][hellblau]{\textcolor{#1}{#2}}
\newcommand{\lila}[2][lila]{\textcolor{#1}{#2}}
\newcommand{\dunkelrot}[2][dunkelrot]{\textcolor{#1}{#2}}
\newcommand{\dunkelgelb}[2][dunkelgelb]{\textcolor{#1}{#2}}


\newcommand{\norm}[1]{\left\lVert#1\right\lVert}
\newcommand{\dv}{\thinspace \mathrm{d}v}
\newcommand{\dw}{\thinspace \mathrm{d}w}
\newcommand{\dx}{\thinspace \mathrm{d}x}
\newcommand{\dy}{\thinspace \mathrm{d}y}
\newcommand{\dt}{\thinspace \mathrm{d}t}
\newcommand{\dr}{\thinspace \mathrm{d}r}
\newcommand{\ds}{\thinspace \mathrm{d}s}
\newcommand{\du}{\thinspace \mathrm{d}u}
\newcommand{\dz}{\thinspace \mathrm{d}z}
\newcommand{\dH}{\thinspace \mathrm{d}H}
\newcommand{\dS}{\thinspace \mathrm{d}S}
\newcommand{\dW}{\thinspace \mathrm{d}W}
\newcommand{\dX}{\thinspace \mathrm{d}X}
\newcommand{\dY}{\thinspace \mathrm{d}Y}
\newcommand{\dZ}{\thinspace \mathrm{d}Z}
\newcommand{\dP}{\thinspace \mathrm{d}P}
\DeclareMathOperator*{\esssup}{ess\,sup}
\DeclareMathOperator*{\argmin}{arg\,min}
%% adjust title page to be centered
% \typearea[18mm]{11}
\usepackage[backend=biber,style=numeric]{biblatex}
\addbibresource{references.bib}
\usepackage{listings}
\definecolor{codebg}{rgb}{0.92,0.92,0.92}
\definecolor{keywordcolor}{rgb}{0.2,0.2,0.7}
\definecolor{stringcolor}{rgb}{0.8,0.1,0.1}
\definecolor{commentcolor}{rgb}{0.4,0.4,0.4}


\lstset{
  backgroundcolor=\color{codebg},
  %frame=single,
  %numbers=left,
  %numberstyle=\tiny\color{gray},
  basicstyle=\ttfamily\small,
  keywordstyle=\color{keywordcolor}\bfseries,
  stringstyle=\color{stringcolor},
  commentstyle=\color{gruen}\itshape,
  language=Python,
  tabsize=2,
  showstringspaces=false
}

\begin{document}

\section{Foundations of Reinforcement Learning}

Wir folgen der Notation aus \cite{Plaat_RL}.
% Grundbegriffe R, MDP, Policy, Value-Funktion
\begin{definition}
    Wir definieren für den Markov decision process (MDP) mit 
    Zustandsraum $\mathcal{S}$, 
    Aktionsraum $\mathcal{A}$, 
    Übergangswahrscheinlichkeiten $P:\mathcal{S} \times \mathcal{A} \times \mathcal{S} \to [0,1]$, 
    Belohnungsfunktion $R:\mathcal{S} \times \mathcal{A} \times \mathcal{S} \to \mathbb{R}$ 
    und Diskontfaktor $\gamma \in [0,1)$ die optimale Value-Funktion als
    %
\end{definition}

% Definition der Value-Funktion
\begin{definition}[Value-Funktion]
    Wir definieren die Value-Funktion für eine Policy $\pi$ als
    \begin{align*}
        V^{\pi}(\bm{s}) 
        = \mathsf{E} 
        \left[
            \sum_{i=0}^{\infty} \gamma^i \cdot r_{t+i+1} \mid \bm{s}_t = \bm{s}
        \right]
    \end{align*}
\end{definition}

% Herleitung V(s) = E(r_t+1 + gamma V(s_t+1) | s_t = s)
\begin{align*}
    V^{\pi}(\bm{s}) 
    &= \mathsf{E} 
    \left[
        \sum_{i=0}^{\infty} \gamma^i \cdot r_{t+i+1} \mid \bm{s}_t = \bm{s}
    \right]
    \\
    &= 
    \mathsf{E} 
    \left[
        r_{t+1} \mid \bm{s}_t = \bm{s}
    \right]
    +
    \rot{\gamma} \cdot
    \mathsf{E} 
    \left[
        \sum_{i=0}^{\infty} \gamma^{i} \cdot r_{t+i+\rot{2}} \mid \bm{s}_t = \bm{s}
    \right]
    \\
    &=
    \mathsf{E} 
    \left[
        r_{t+1} \mid \bm{s}_t = \bm{s}
    \right]
    +
    \rot{\gamma} \cdot
    \mathsf{E} 
    \left[
        \underbrace{        
            \mathsf{E}
            \left[
                \sum_{i=0}^{\infty} \gamma^{i} \cdot r_{t+i+\rot{2}} \mid \bm{s}_{t+1}
            \right]
            }_{
                V^{\pi}(\bm{s}_{t+1})
            }
        \mid \bm{s}_t = \bm{s}
    \right]
    \\
    \Rightarrow
    V^{\pi}(\bm{s})
    &=
    \mathsf{E} 
    \left[
        r_{t+1} + \gamma \cdot V^{\pi}(\bm{s}_{t+1}) \mid \bm{s}_t = \bm{s}
    \right] 
\end{align*}

Denn es gilt hier für die $\sigma$-Algebren, 
dass $\sigma(\bm{s}_t ) \subseteq \sigma(\bm{s}_{t+1})$ gilt. 
Explizit ausgeschrieben in Erwartungswerten ergibt sich:

% Herleitung V(s) = E_a~pi [ E_s'~T_a [ R_a(s,s') + gamma V(s') ] ]
\begin{align*}
    V^{\pi}(\bm{s})
    &= \mathsf{E} 
        \left[ 
            r_{t+1} + \gamma V^{\pi}(\bm{s}_{t+1}) \mid \bm{s}_t = \bm{s} 
        \right] \\
    &= \mathsf{E}_{a_t \sim \pi(\cdot|\bm{s})} 
        \left[ 
            \vphantom{\int}
            \mathsf{E} 
            \left[ 
                r_{t+1} + \gamma V^{\pi}(\bm{s}_{t+1}) \mid \bm{s}_t = \bm{s}, a_t = a 
            \right] 
        \right] \\
    &= \mathsf{E}_{a_t \sim \pi(\cdot|\bm{s})} 
    \left[ 
        \vphantom{\int}
        \mathsf{E}_{\bm{s}_{t+1} \sim T_a(\bm{s})} 
        \left[ 
            R_{a_t}(\bm{s}_t, \bm{s}_{t+1}) + \gamma V^{\pi}(\bm{s}_{t+1}) 
        \right] 
    \right] \\
    &= \mathsf{E}_{a \sim \pi(\cdot|\bm{s})} 
        \left[ 
            \vphantom{\int}
            \mathsf{E}_{\bm{s}' \sim T_a(\bm{s})} 
            \left[ 
                R_{a}(\bm{s}, \bm{s}') + \gamma V^{\pi}(\bm{s}') 
            \right] 
        \right]
\end{align*}

% Definition der Q-Funktion
\begin{definition}[Q-Function]
    Wir definieren die Q-Function als
    \begin{align*}
        Q^{\pi}(\bm{s}, \gruen{a}) 
        =
        \mathsf{E}
        \left[
            \vphantom{\int}
            \sum_{i=0}^{\infty} \gamma^i \cdot r_{t+i+1} 
            \mid \bm{s}_t = \bm{s}, \gruen{a_t = a}
        \right]
        \Rightarrow
        V^{\pi}(\bm{s})
        =
        \mathsf{E}_{\gruen{a} \sim \pi(\cdot|\bm{s})}
        \left[
            Q^{\pi}(\bm{s}, \gruen{a})
        \right]
    \end{align*}    
\end{definition}

Herleitung der Bellmann Gleichung für die Q-Funktion:
% Q(s,a) = E_s'~T_a [ R_a(s,s') + gamma E_a'~pi [ Q(s',a') ] ]
% Q(s,a) = E_s'~T_a [ R_a(s,s') + gamma V(s') ]
\begin{align*}
    Q^{\pi}(\bm{s}, a) 
    &= 
    \mathsf{E}
    \left[
        \vphantom{\int}
        \sum_{i=0}^{\infty} \gamma^i \cdot r_{t+i+1} 
        \mid \bm{s}_t = \bm{s}, a_t = a
    \right] \\
    &=
    \mathsf{E}
    \left[  
        \mathsf{E}
        \left[
            r_{t+1} + \sum_{i=1}^{\infty} \gamma^i \cdot r_{t+i+1} 
            \mid \bm{s}_t = \bm{s}, a_t = a, \bm{s}_{t+1} = \bm{s}'
            \right]
            \mid \bm{s}_t = \bm{s}, a_t = a
        \right] \\
        &= 
        \mathsf{E}
        \left[
            r_{t+1} + \gamma \cdot 
            \underbrace{
                \mathsf{E}
                \left[
                    \sum_{i=0}^{\infty} \gamma^{i} 
                    \cdot r_{t+i+2} \mid \bm{s}_{t+1} = \bm{s}'
                \right]
            }_{
                V^{\pi}(\bm{s}_{t+1})
            } 
            \mid \bm{s}_t = \bm{s}, a_t = a
        \right] \\
        &= 
        \mathsf{E}_{\bm{s}' \sim T_a(\bm{s})}
        \left[
            \vphantom{\int}
            R_a(\bm{s}, \bm{s}') + \gamma V^{\pi}(\bm{s}') 
        \right]
        \\
        &= 
        \mathsf{E}_{\bm{s}' \sim T_a(\bm{s})}
        \left[
            \vphantom{\int}
            R_a(\bm{s}, \bm{s}') + \gamma 
            \mathsf{E}_{a' \sim \pi(\cdot|\bm{s}')}
            \left[
                \vphantom{\int}
                Q^{\pi}(\bm{s}', a')
            \right]
        \right]
\end{align*}

    Wobei wir im letzten Schritt die folgende Beziehung genutzt haben:
    % V(s) = E_a [Q(s,a)]
    \begin{align*}
        V^{\pi}(\bm{s})
        &=
        \mathsf{E}_{a \sim \pi(\cdot|\bm{s})}
        \left[
            Q^{\pi}(\bm{s}, a)
        \right]
    \end{align*}

\subsection{Bellman Optimality Equations}

% Bellmann Optimality Equations 
% V* = max_a E_s'~T_a [ R_a(s,s') + gamma V*(s') ]
\begin{satz}[Optimality Equation für $V$] \label{Bellman_Optimality_V}
    Die optimale Value-Funktion $V^*(\bm{s})$ erfüllt:
    \begin{align*}
        V^{*}(\bm{s})
        =
        \max_{a} \mathsf{E}_{\bm{s}' \sim T_a(\bm{s})}
        \left[
            \vphantom{\int}
            R_{a}(\bm{s}, \bm{s}') + \gamma V^{*}(\bm{s}')
        \right]
        \quad
        \pi^*(\bm{s}) 
        =
        \arg\max_{a} \mathsf{E}_{\bm{s}' \sim T_a(\bm{s})}
        \left[
            \vphantom{\int}
            R_{a}(\bm{s}, \bm{s}') + \gamma V^{*}(\bm{s}')
        \right]
    \end{align*}
\end{satz}

\begin{beweis}
  Wir setzen $V^* (\bm{s}) = \max_{\pi} V^{\pi}(\bm{s})$. 
  Für jede Policy $\pi$ gilt nun $V^*(\bm{s}') \geq V^{\pi}(\bm{s}')$ und somit:
  % Herleitung Bellman Optimality Equation
  % V*(s) = max_pi V^pi(s) = max_a E_s'~T_a [ R_a(s,s') + gamma V*(s') ]
  \begin{align*}
      V^{*}(\bm{s})
      &= \max_{\pi} 
      \mathsf{E}_{a \sim \pi(\cdot|\bm{s})} 
      \left[ 
          \vphantom{\int}
          \mathsf{E}_{\bm{s}' \sim T_a(\bm{s})} 
          \left[ 
              R_{a}(\bm{s}, \bm{s}') + \gamma \rot{V^{\pi}}(\bm{s}')
          \right] 
      \right] 
      \rot{\leq}\;
      \max_{\pi}
      \mathsf{E}_{a \sim \pi(\cdot|\bm{s})}
      \left[ 
          \vphantom{\int}
          \mathsf{E}_{\bm{s}' \sim T_a(\bm{s})} 
          \left[
              R_{a}(\bm{s}, \bm{s}') + \gamma V^{*}(\bm{s}')
          \right]
      \right]
  \end{align*}

  Das Maximum über alle Policies $\pi$ wird erreicht durch die Policy,
  die in jedem Zustand $\bm{s}$ die Aktion $a$ auswählt, die das Maximum 
  in der inneren Erwartung erreicht. Somit gilt:
  %
  \begin{align*}
      V^{*}(\bm{s})
      &\leq \max_{a} \mathsf{E}_{\bm{s}' \sim T_a(\bm{s})}
      \left[
          \vphantom{\int}
          R_{a}(\bm{s}, \bm{s}') + \gamma V^{*}(\bm{s}')
      \right]
  \end{align*}

  Wir konstruieren nun die greedy Policy $\pi^*$ mittels:
  % Definition pi* = argmax_a E_s'~T_a [ R_a(s,s') + gamma V*(s') ]
  %
  \begin{align*}
      \pi^*(\bm{s}) 
      =
      \arg\max_{a} \mathsf{E}_{\bm{s}' \sim T_a(\bm{s})}
      \left[
          \vphantom{\int}
          R_{a}(\bm{s}, \bm{s}') + \gamma V^{*}(\bm{s}')
      \right]
  \end{align*}

  Für diese Policy $\pi^*$ gilt nun:
  % Herleitung V^{pi*}(s) = max_a E_s'~T_a [ R_a(s,s') + gamma V*(s') ]
  %
  \begin{align*}
      V^{\pi^*}(\bm{s})
      &= 
      \max_{a} \mathsf{E}_{\bm{s}' \sim T_a(\bm{s})}
      \left[
          \vphantom{\int}
          R_{a}(\bm{s}, \bm{s}') + \gamma V^{*}(\bm{s}')
      \right]
      \geq V^{*}(\bm{s})
      \Rightarrow
      V^{\pi^*}(\bm{s}) = V^{*}(\bm{s})
  \end{align*}
\end{beweis}

% Bellmann Optimality Equations for Q-function
% Q*(s,a) = E_s'~T_a [ R_a(s,s') + gamma max_a' Q*(s',a') ]
\begin{satz}[Optimality Equation für Q-Funktion] \label{Bellman_Optimality_Q}
    Die optimale Q-Funktion $Q^*(\bm{s}, a)$ erfüllt:
    %
    \begin{align*}
        Q^{*}(\bm{s}, a) =
        \mathsf{E}_{\bm{s}' \sim T_a(\bm{s})}
        \left[
            \vphantom{\int}
            R_a(\bm{s}, \bm{s}') + \gamma 
            \max_{a'} Q^{*}(\bm{s}', a')
        \right]
        \quad
        \pi^{*}(\bm{s}) = \arg\max_{a} Q^{*}(\bm{s}, a)
    \end{align*}
\end{satz}

\begin{beweis}
  Für die optimale Q-Funktion gilt somit
  % Q* = max_pi Q(a,s)
  \begin{align*}
      Q^{*}(\bm{s}, a) 
      =
      \max_{\pi} Q^{\pi}(\bm{s}, a)
      \Rightarrow
      Q^{*}(\bm{s}, a) =
      \mathsf{E}_{\bm{s}' \sim T_a(\bm{s})}
      \left[
          \vphantom{\int}
          R_a(\bm{s}, \bm{s}') + \gamma 
          \max_{a'} Q^{*}(\bm{s}', a')
      \right]
  \end{align*}

  Für die optimale Policy gilt:
  \begin{align*}
      \pi^{*}(\bm{s}) = a^* = \arg\max_{a} Q^{*}(\bm{s}, a)
  \end{align*}
\end{beweis}

\subsection{Value Iteration}

Wir definieren die Abbildung $\mathcal{T}$ als
\begin{align*}
    \mathcal{T}[V](\bm{s}, a)
    &:=
    \max_{a}
    \mathsf{E}_{\bm{s}' \sim T_a(\bm{s})}
    \left[
        \vphantom{\int}
        R_a(\bm{s}, \bm{s}') + \gamma 
        V(\bm{s}', a')
    \right]
\end{align*}

Die Bellmann Optimalitätsgleichung \ref{Bellman_Optimality_V} sagt nun aus, 
dass $V^*$ der Fixpunkt der Abbildung $\mathcal{T}$ ist, 
also $\mathcal{T}[V^*] = V^*$ gilt. 
Die Fixpunktiteration $V_{k+1} = \mathcal{T}[V_k]$ konvergiert somit gegen $V^*$.
%

\begin{beweis}
  Wir zeigen zunächst, dass $\mathcal{T}$ eine Kontraktion ist:
  % T(V) - T(U) <= gamma ||V-U||_infty
  \begin{align*}
      \left| \mathcal{T}[V] - \mathcal{T}[U] \right|
      &=
      \left|
          \lila{\max_{a}}
          \mathsf{E}_{\bm{s}' \sim T_a(\bm{s})}
          \left[
              \vphantom{\int}
              R_a(\bm{s}, \bm{s}') + \gamma 
              V(\bm{s}')
          \right]
          -
          \lila{\max_{a}}
          \mathsf{E}_{\bm{s}' \sim T_a(\bm{s})}
          \left[
              \vphantom{\int}
              R_a(\bm{s}, \bm{s}') + \gamma 
              U(\bm{s}')
          \right]
      \right| \\
      &\leq
      \lila{\max_{a}}
      \left|
          \mathsf{E}_{\bm{s}' \sim T_a(\bm{s})}
          \left[
              \vphantom{\int}
              R_a(\bm{s}, \bm{s}') + \gamma 
              V(\bm{s}')
          \right]
          -
          \mathsf{E}_{\bm{s}' \sim T_a(\bm{s})}
          \left[
              \vphantom{\int}
              R_a(\bm{s}, \bm{s}') + \gamma 
              U(\bm{s}')
          \right]
      \right|
      \\
      &=
      \max_{a}
      \left|
          \mathsf{E}_{\bm{s}' \sim T_a(\bm{s})}
          \left[
              \vphantom{\int}
              \gamma
              (V(\bm{s}') - U(\bm{s}'))
          \right]
      \right|
      \\
      &\leq
      \max_{a}
      \mathsf{E}_{\bm{s}' \sim T_a(\bm{s})}
      \left[
          \vphantom{\int}
          \gamma
          \cdot
          \underbrace{
              \rot{\sup_{\bm{s}'} |} V(\bm{s}') - U(\bm{s}') \rot{|}
          }_{
              \|V - U\|_{\infty}
          }
      \right]
      \\
      &=
      \gamma \|V - U\|_{\infty}
  \end{align*}

  Nach dem Fixpunktsatz von Banach \cite{sutton_barto_2018}
  konvergiert die Iteration $V_{k+1} = \mathcal{T}(V_k)$ gegen $V^*$ für $k \to \infty$.
\end{beweis}

\begin{algorithm}
    \caption{Value Iteration}
    \begin{algorithmic}
        \State Initialize $V_0(s)$ arbitrarily for all $s \in \mathcal{S}$
        \State $k \gets 0$
        \Repeat
        \For{each $s \in \mathcal{S}$}
        \State $V_{k+1}(s) \gets \max_{a} 
        \underbrace{
            \mathsf{E}_{s' \sim T_a(s)} 
            \left[ 
                \vphantom{\int}
                R_a(s, s') + \gamma V_k(s') 
            \right]}_{Q_{k}(s,a)}
            =
            \max_a \sum_{s'} T_a(s,s') 
            \left[ 
                \vphantom{\int}
                R_a(s,s') + \gamma V_k(s') 
            \right]
        $
        \EndFor
        \State $k \gets k + 1$
        \Until{$|V_k - V_{k-1}| < \varepsilon$}
        \State \Return $V_k$
    \end{algorithmic}
\end{algorithm}

\newpage

\subsection{Q-Learning}

Wir betrachten die Optimalitätsgleichung für die Q-Funktion \ref{Bellman_Optimality_Q}:
%
\begin{align*}
    Q^{*}(\bm{s}, a) =
    \mathsf{E}_{\bm{s}' \sim T_a(\bm{s})}
    \left[
        \vphantom{\int}
        R_a(\bm{s}, \bm{s}') + \gamma 
        \max_{a'} Q^{*}(\bm{s}', a')
    \right]
\end{align*}

und definieren erneut den Operator $\mathcal{T}$ als
%
\begin{align*}
    \mathcal{T}[Q](\bm{s}, a)
    &:=
    \mathsf{E}_{\bm{s}' \sim T_a(\bm{s})}
    \left[
        \vphantom{\int}
        R_a(\bm{s}, \bm{s}') + \gamma 
        \max_{a'} Q(\bm{s}', a')
    \right]
\end{align*}

und erhalten erneut die Fixpunktgleichung $\mathcal{T}[Q^*] = Q^*$. Die Fixpunktiteration
$Q_{k+1} = \mathcal{T}[Q_k]$ konvergiert somit gegen $Q^*$, 
allerdings ist die Berechnung der Erwartungswerte nicht möglich, 
da wir das Modell $\bm{s}' = T_a(\bm{s})$ nicht kennen.
%

% Hier konkreter Q-Learning Algorithmus
%--------------------------------------------------------------
\begin{algorithm}
    \caption{Q-Learning}
    \begin{algorithmic}
        \State Given Environment, 
        learning rate $\alpha \in (0,1]$, 
        discount factor $\gamma \in [0,1)$,
        exploration rate $\epsilon \in (0,1)$
        \State Initialize $Q(s,a)$ arbitrarily for all $s \in \mathcal{S}, a \in \mathcal{A}$
        \State $Q(\text{terminal}, \cdot) = 0$ \Comment{\lila{Terminale Zustände haben Wert 0}}
        \For{each episode}
            \State Initialize state $s$
            \While{state $s$ is not terminal}
                \State Choose action $a = \epsilon$-greedy$(Q[s,\cdot],\epsilon)$ 
                \Comment{\gruen{Exploration vs. Exploitation}}
                \State $(s',r) = Environment(s,a)$
                \State $Q[s,a] \leftarrow (1-\alpha) \cdot Q[s,a] 
                + \alpha \cdot (r + \gamma \cdot \max_{a'} Q[s',a'])$
                \Comment{\rot{Bootstrapping und TD-Learning}}
                \State $s \leftarrow s'$
            \EndWhile
        \EndFor
    \end{algorithmic}
\end{algorithm}

\begin{beweis}
  Wir prüfen die Vorraussetzungen für die Konvergenz von 
  Q-Learning aus \cite{tsitsiklis1992asynchronous} und folgern die 
  Konvergenz gegen $Q^*$ fast sicher.
  %
  \begin{itemize}
      \item Der Beweis, dass $\mathcal{T}$ eine Kontraktion ist,
      verläuft analog zum Beweis für die Value-Funktion.
  %
      \begin{align*}
          \left| \mathcal{T}[Q] - \mathcal{T}[U] \right|
          &=
          \left|
              \mathsf{E}_{\bm{s}' \sim T_a(\bm{s})}
              \left[
                  \vphantom{\int}
                  R_a(\bm{s}, \bm{s}') + \gamma 
                  \max_{a'} Q(\bm{s}', a')
              \right]
              -
              \mathsf{E}_{\bm{s}' \sim T_a(\bm{s})}
              \left[
                  \vphantom{\int}
                  R_a(\bm{s}, \bm{s}') + \gamma 
                  \max_{a'} U(\bm{s}', a')
              \right]
          \right| \\
          &=
          \left|
              \mathsf{E}_{\bm{s}' \sim T_a(\bm{s})}
              \left[
                  \vphantom{\int}
                  \gamma
                  \left(
                      \max_{a'} Q(\bm{s}', a') - \max_{a'} U(\bm{s}', a')
                  \right)
              \right]
          \right| \\
          &\leq
          \gamma
          \mathsf{E}_{\bm{s}' \sim T_a(\bm{s})}
          \left[
              \vphantom{\int}
              \gamma
              \cdot
              \underbrace{
                  \rot{\max_{a'} |} Q(\bm{s}', a') - U(\bm{s}', a') \rot{|}
              }_{
                  \leq \|Q - U\|_{\infty}
              }
          \right]
          \\
          &\leq
          \gamma \|Q - U\|_{\infty}
      \end{align*}
      %
      \item Wir verwenden folgende Update-Regel, um die Q-Funktion zu approximieren:
  %
      \begin{align*}
          Q_{k+1}(\bm{s}, a) 
          &=  Q_k(\bm{s}, a)
          + \alpha \cdot 
          \left[
              \underbrace{
                  R_a(\bm{s},\bm{s}') + \gamma \cdot \max_{a'} Q_k(\bm{s}', a')
                  }_{\text{Sample of }T[Q_k]} 
              - Q_k(\bm{s}, a)
          \right]
          \\
          &=
          Q_k(\bm{s}, a)
          + \alpha \cdot
          \left[
              \rot{\mathcal{T}[Q_k](\bm{s}, a)} - Q_k(\bm{s}, a) +
              \underbrace{
                  R_a(\bm{s},\bm{s}') + \gamma \cdot \max_{a'} Q_k(\bm{s}', a')
                  \rot{- \mathcal{T}[Q_k](\bm{s}, a)}}_{
                      \eta_k
                  }
          \right]
          \\
          &=  Q_k(\bm{s}, a)
          + \alpha \cdot
          \left[
              \vphantom{\int}
              \mathcal{T}[Q_k](\bm{s}, a) - Q_k(\bm{s}, a)
              + \eta_k
          \right]
      \end{align*}
      %
      Für den Erwartungswert des Rauschterms $\eta_k$ gilt:
      %
      \begin{align*}
          \mathsf{E}_{\bm{s}' \sim T_a(\bm{s})}
          \left[
              \eta_k \mid Q_k
          \right]
          &=
          \mathsf{E}_{\bm{s}' \sim T_a(\bm{s})}
          \left[
              R_a(\bm{s},\bm{s}') + \gamma \cdot \max_{a'} Q_k(\bm{s}', a')
              - 
              \underbrace{\mathcal{T}[Q_k](\bm{s}, a)}_{
                  \text{deterministic}
              }
              \mid Q_k
          \right]
          = 0
      \end{align*}
      %
      \item Wir definieren die $\epsilon$-Greedy-Policy:
  %
      \begin{align*}
          \pi_{\epsilon}(a, \bm{s})
          &:=
          \begin{cases}
              1- \epsilon + \frac{\varepsilon}{|\mathcal{A}|} & \text{if } a = \arg\max_{a'} Q[\bm{s}, a'] \\
              \frac{\varepsilon}{|\mathcal{A}|} & \text{else}
          \end{cases}
          =
          (1-\epsilon) \cdot \mathds{1}_{a = \arg\max_{a'} Q[\bm{s}, a']} 
          + \frac{\epsilon}{|A|}
      \end{align*}
  \end{itemize}
\end{beweis}


\newpage

\printbibliography

\end{document}